\chapter{Results and Operational Robustness Analysis}
\label{chap:results}

Il presente capitolo espone i risultati sperimentali della pipeline di valutazione
\textsc{FAIR-Lab}, con particolare attenzione alla robustezza operativa di sistemi
di classificazione automatizzata di immagini basati su \gls{ai} in uno scenario di
triage forense digitale. L'obiettivo dell'analisi non è ordinare i modelli
esclusivamente in base alla loro accuratezza nominale, bensì valutare se gli
output prodotti rimangano affidabili in presenza di input puliti, campioni
out-of-distribution, trasformazioni anti-forensi e perturbazioni adversariali. In
questa prospettiva, assumono particolare rilievo i falsi negativi relativi alla
classe \texttt{weapon}, le predizioni ad alta confidenza su campioni fuori
distribuzione e i requisiti di tracciabilità necessari per confrontare modelli
proxy trasparenti con strumenti forensi commerciali operanti come sistemi
black-box.

Nel corso del capitolo viene mantenuta una distinzione metodologica esplicita
tra i modelli proxy trasparenti, per i quali risultano disponibili architetture,
checkpoint, rappresentazioni intermedie e output predittivi, e gli strumenti
forensi commerciali, che operano come sistemi black-box o quasi black-box, poiché
la loro logica interna di classificazione è proprietaria e non direttamente
ispezionabile. Tale distinzione non è soltanto tecnica, ma incide direttamente
sulla tracciabilità, sull'auditabilità e sull'interpretazione probatoria degli
output assistiti da \gls{ai} all'interno di workflow forensi.

L'organizzazione del capitolo segue la frequenza operativa attesa delle diverse
condizioni sperimentali in un contesto investigativo. Le prestazioni su input
puliti definiscono la baseline di riferimento; il comportamento su campioni
out-of-distribution e su trasformazioni anti-forensi rappresenta condizioni
realistiche che possono emergere nel corso di indagini digitali; le perturbazioni
adversariali descrivono invece tentativi deliberati di evasione e costituiscono
la condizione di stress più severa considerata in questa valutazione. La
valutazione degli strumenti forensi commerciali è trattata come livello
black-box distinto, introdotto dopo la presentazione dei risultati controllati dei
modelli proxy e dopo la definizione del forensic evaluation bundle.

% ─────────────────────────────────────────────────────────────────────────────
\section{Experimental Evaluation Status and Dataset Overview}
\label{sec:results-status}
% ─────────────────────────────────────────────────────────────────────────────

La valutazione sperimentale si basa sugli artefatti generati attraverso la
metodologia descritta nel Capitolo~\ref{chap:methodology}. Allo stato attuale,
risultano completate la costruzione del dataset, il congelamento semantico delle
classi, la generazione degli split, l'addestramento dei modelli proxy, la
generazione delle perturbazioni, la valutazione dei modelli proxy e la
costruzione del forensic evaluation bundle. La valutazione black-box degli
strumenti forensi commerciali è invece considerata una fase operativa separata,
poiché dipende dalla disponibilità dei tool, dai moduli effettivamente abilitati,
dai formati di esportazione e dalla successiva normalizzazione degli output
proprietari.

Questa separazione è metodologicamente rilevante. Il livello dei modelli proxy
fornisce una baseline trasparente e riproducibile, poiché architetture,
checkpoint, fold di addestramento, predizioni e metriche sono disponibili per
ispezione e verifica. Il livello degli strumenti forensi commerciali riflette
invece una condizione operativa differente: l'analista interagisce con strumenti
professionali la cui logica interna di classificazione non è pienamente
accessibile, e la valutazione deve quindi basarsi sugli output osservabili, sulle
categorie esportate, sui tag assegnati, sugli eventuali valori di confidenza e
sul matching hash-based rispetto alla ground truth sperimentale.

La Tabella~\ref{tab:results-evaluation-status} riporta lo stato corrente delle
principali componenti sperimentali, fornendo una sintesi audit-ready del workflow
di valutazione al momento della redazione.

\begin{table}[H]
\centering
\caption{Stato corrente del workflow di valutazione sperimentale.}
\label{tab:results-evaluation-status}
\begin{tabularx}{\textwidth}{@{}l l X@{}}
\toprule
\textbf{Componente} & \textbf{Stato} & \textbf{Note} \\
\midrule

Dataset finale
&
\textbf{Congelato}
&
1500 immagini suddivise in tre classi: 500 \texttt{weapon}, 500
\texttt{non\_weapon} e 500 \texttt{ood}. \\

Subset binario
&
\textbf{Congelato}
&
1000 immagini, composte da 500 campioni \texttt{weapon} e 500 campioni
\texttt{non\_weapon}. Il subset è utilizzato per gli split clean, per gli
attacchi adversariali e per le trasformazioni anti-forensi. \\

Split clean e OOD
&
\textbf{Completati}
&
Sono stati generati cinque fold binari clean e un insieme di valutazione
out-of-distribution separato, entrambi derivati dai manifest congelati. \\

Addestramento dei modelli proxy
&
\textbf{Completato}
&
Sono stati prodotti checkpoint fold-aware per \texttt{efficientnet\_b0},
\texttt{resnet18} e \texttt{clip}. \\

Valutazione dei modelli proxy
&
\textbf{Completata}
&
Il livello proxy è stato valutato su input clean, OOD, adversariali e
anti-forensi, producendo output a livello di predizione e tabelle metriche
finali. \\

Attacchi adversariali
&
\textbf{Generati}
&
Sono state generate cinque famiglie di perturbazioni: \texttt{fgsm},
\texttt{superdeepfool}, \texttt{sigma\_zero}, \texttt{one\_pixel} e
\texttt{color\_shift}. \\

Trasformazioni anti-forensi
&
\textbf{Generate}
&
Sono state generate cinque trasformazioni model-agnostic:
\texttt{jpeg\_recompression}, \texttt{resample\_resize},
\texttt{gaussian\_blur}, \texttt{histogram\_modification} e
\texttt{contrast\_stretching}. \\

Forensic evaluation bundle
&
\textbf{Generato}
&
11500 file predisposti per la valutazione black-box degli strumenti forensi,
includendo campioni clean, OOD, adversariali e anti-forensi. \\

Magnet AXIOM / Magnet.AI
&
\textbf{In attesa}
&
Bundle sottoposto o in fase di importazione; la valutazione black-box e la
normalizzazione degli export non sono ancora consolidate nelle tabelle metriche
finali. \\

X-Ways Forensics / Excire
&
\textbf{In attesa}
&
La valutazione dipende dalla disponibilità del tool, dai moduli di
classificazione visuale abilitati e dal formato di esportazione disponibile. \\

Cellebrite UFED
&
\textbf{In attesa}
&
La valutazione dipende dalla disponibilità del tool e dal supporto
all'importazione e all'esportazione di risultati di classificazione delle
immagini. \\

Oxygen Forensic Detective
&
\textbf{In attesa}
&
La valutazione dipende dalla disponibilità del tool, dai moduli \gls{ai}
abilitati e dagli output di classificazione effettivamente esportabili. \\

Casi di studio di explainability
&
\textbf{Pianificati}
&
Le attribution map basate su Integrated Gradients saranno calcolate su casi di
fallimento rappresentativi dei modelli proxy; tali risultati avranno valore
qualitativo e diagnostico, senza costituire metriche primarie di valutazione. \\

\bottomrule
\end{tabularx}
\end{table}

Il corpus sperimentale risultante è organizzato in condizioni clean,
out-of-distribution, adversariali e anti-forensi. I campioni binari clean e
perturbati derivano dal subset ufficiale \texttt{weapon}/\texttt{non\_weapon},
mentre i campioni OOD sono valutati separatamente e non sono utilizzati come
target degli attacchi. Tale distinzione preserva la differenza tra classificazione
binaria nominale, comportamento fuori distribuzione e robustezza rispetto a input
manipolati.

La Tabella~\ref{tab:results-corpus-composition} riporta la composizione
consolidata del corpus di valutazione utilizzato nei livelli proxy-model e
forensic-tool.

\begin{table}[H]
\centering
\caption{Composizione del corpus sperimentale e del forensic evaluation bundle.}
\label{tab:results-corpus-composition}
\begin{tabularx}{0.88\textwidth}{@{}l r X@{}}
\toprule
\textbf{Condizione} & \textbf{File} & \textbf{Ruolo nella valutazione} \\
\midrule

Clean
&
1000
&
Valutazione binaria nominale sui campioni \texttt{weapon} e
\texttt{non\_weapon}. \\

OOD
&
500
&
Insieme di valutazione out-of-distribution e borderline, mantenuto separato dal
task di classificazione binaria. \\

Adversarial
&
5000
&
Cinque famiglie di perturbazioni applicate al subset binario di 1000 immagini,
con 1000 file generati per ciascuna famiglia di perturbazione. \\

Anti-forensic
&
5000
&
Cinque trasformazioni model-agnostic applicate al subset binario di 1000
immagini, con 1000 file generati per ciascuna trasformazione. \\

\midrule
\textbf{Totale}
&
\textbf{11500}
&
Forensic evaluation bundle completo, predisposto per la valutazione dei modelli
proxy e per il processamento mediante strumenti forensi black-box. \\

\bottomrule
\end{tabularx}
\end{table}

Il forensic evaluation bundle costituisce il collegamento operativo tra il
livello sperimentale controllato e il livello di valutazione basato su strumenti
forensi professionali. Esso fornisce una directory di input cieca e
semanticamente neutra per i tool forensi, preservando al contempo ground truth,
informazioni di provenienza, metadati sulle perturbazioni e mapping hash-based in
file separati. In termini operativi, gli strumenti forensi devono processare
soltanto la directory cieca di input, mentre il livello dei metadati è riservato
alle fasi successive di matching, normalizzazione e audit degli export.

Questo disegno riduce il rischio di bias indotto dai path o dall'analista, poiché
i tool non sono esposti a label di classe, nomi degli attacchi, identificativi di
sorgente o informazioni sui fold durante l'importazione. Al tempo stesso, la
presenza di mapping SHA256 e MD5 consente di ricondurre ciascun risultato
esportato al corrispondente artefatto sperimentale. Tale requisito di
tracciabilità è centrale per l'interpretazione forense dei risultati, poiché
permette di distinguere tra errori di classificazione, file non supportati,
mancate rilevazioni, output non categorizzati e limitazioni specifiche del formato
di export del singolo tool.

Le sezioni successive presentano i risultati quantitativi, a partire dalla
baseline clean descritta nella Sezione~\ref{sec:results-clean-baseline}.

\section{Clean Baseline Performance}
\label{sec:results-clean-baseline}
% T1: Accuracy, balanced accuracy, precision, recall, F1,
% false positive rate, false negative rate per modello
% F1: Confusion matrix per modello

\section{Out-of-Distribution Behavior and False Positive Risk}
\label{sec:results-ood}
% T2: weapon rate, confidence media, high-confidence rate per modello
% F2: distribuzione confidence su OOD, se disponibile

\section{Robustness Against Anti-Forensic Transformations}
\label{sec:results-antiforensic}
% T3: accuracy drop per modello e trasformazione
% F3: bar chart o heatmap drop anti-forensic

\section{Robustness Against Adversarial Perturbations}
\label{sec:results-adversarial}
% T4: accuracy drop per modello e attacco
% F4: heatmap drop adversarial
% Optional: curva accuracy vs epsilon solo se viene eseguito uno sweep FGSM

\section{Forensic Evaluation Bundle Composition and Black-Box Readiness}
\label{sec:results-bundle}
% T5: composizione bundle da 11500 file
% Protocollo di matching e normalizzazione degli output dei tool commerciali

\section{Commercial Forensic Tools Evaluation}
\label{sec:results-commercial}
% T6: matrice risultati/tool availability con celle pending dichiarate
% Formula trasparente per risultati non ancora disponibili

\section{Comparative Robustness Discussion}
\label{sec:results-comparative}
% T7: sintesi comparativa dei proxy model e, quando disponibili,
% dei tool commerciali normalizzati

\section{Explainability Analysis and Representative Failure Cases}
\label{sec:results-xai}
% F5: Integrated Gradients su casi rappresentativi clean vs perturbed
% Gallery completa → Appendice

\section{Operational Forensic Implications and Limitations}
\label{sec:results-implications}
% Falsi negativi, tracciabilità, affidabilità operativa,
% calibrazione confidence, trasferibilità degli attacchi,
% limiti dei tool black-box